Die Mutter ist die Drehachse. Drückst du senkrecht zum Schlüssel, ist der Hebelarm der Abstand $d$ von der Mutter; aus $M = F\cdot d$ folgt $F = M/d$.

\begin{abcliste}
\abc
Mit dem Schlüssel:
\begin{align}
 F_1 &= \FaF \\
   &= \frac{\M}{\la} \\
   &= \Fa \approx \result{\FaP{0}{2}}
\end{align}
Das ist etwa die Gewichtskraft einer Masse von $F_1/g \approx \mFP{0}{2}$ (mit \gO): mehr, als die meisten Menschen mit den Armen drücken können.

\abc
Mit dem Rohr:
\begin{align}
 F_2 &= \FbF \\
   &= \frac{\M}{\lb} \\
   &= \Fb \approx \result{\FbP{0}{3}}
\end{align}
Ein längerer Hebelarm: dasselbe Drehmoment mit weniger Kraft.
\end{abcliste}
